%! Author = lili
%! Date = 7/9/26

\chapter{Tutorium 18.05.2026}
\untit{Präsenzaufgaben 1.4: Unabhängigkeit \hfill 18.--22.\ Mai 2026}

\section{PÜ 1.4.4.}
Ein (fairer) Würfel wird $20$ Mal geworfen.
\begin{auf}[Geben Sie einen geeigneten Wahrscheinlichkeitsraum $(\eraum,\pe)$ an, um die
folgenden Aufgabenteile bearbeiten zu können.]
    NOCH KEINE LÖSUNG % TODO NOCH KEINE LÖSUNG
\end{auf}

\begin{auf}[Geben Sie die folgenden Ereignisse als Teilmengen von $\eraum$ an, und berechnen
Sie ihre Wahrscheinlichkeiten:]
    \index{Aufgabe!Wahrscheinlichkeiten berechnen}
    \begin{itemize}
        \item $A$: Alle $20$ Würfe haben dasselbe Ergebnis.
        \item $B$: Die größte geworfene Zahl ist größer als $4$.
        \item $C$: Der erste Wurf zeigt eine $4$.
        \item $D$: Die größte geworfene Zahl ist kleiner als $5$.
    \end{itemize}
    NOCH KEINE LÖSUNG % TODO NOCH KEINE LÖSUNG
\end{auf}

\begin{auf}[Fragen Beantworten]
    \index{Aufgabe!Unabhängigkeit bewerten}
    \begin{enumerate}
        \item Sind $A$ und $B$ unabhängig?
        \item Sind $A$ und $C$ unabhängig?
        \item Sind $B$ und $C$ unabhängig?
        \item Sind $A, B, C$ unabhängig?
        \item Sind $A, B, C, D$ unabhängig?
    \end{enumerate}
    NOCH KEINE LÖSUNG % TODO NOCH KEINE LÖSUNG
\end{auf}

\section{PÜ 1.4.5.}

Sei $(X,Y)$ uniform verteilt auf der Menge
\[
    A = \left\{(x,y)\in\rz^2 : 0 \leq x,y \leq \frac{1}{2}\right\}
    \cup
    \left\{(x,y)\in\rz^2 : \frac{1}{2} < x,y \leq 1\right\}.
\]

\begin{auf}[Skizzieren Sie die Menge $A$.]
    NOCH KEINE LÖSUNG % TODO NOCH KEINE LÖSUNG
\end{auf}

\begin{auf}\subsubsection*{Zeigen Sie, dass sowohl $X$ als auch $Y$ uniform verteilt auf $\left[0,1\right]$ ist.}
\index{Aufgabe!uniforme Verteilung zeigen}
    NOCH KEINE LÖSUNG % TODO NOCH KEINE LÖSUNG
\end{auf}

\begin{auf}[Können Sie zwei Mengen $C, D \subseteq \rz$ finden derart, dass die Ereignisse
$\{X \in C\}$ und $\{Y \in D\}$ nicht unabhängig sind?]
    NOCH KEINE LÖSUNG % TODO NOCH KEINE LÖSUNG
\end{auf}